\section{Finite certificates for modules and automorphisms}

The following criteria explain what the finite data in the applications
must certify. They use ordinary linear algebra and group actions. A list
of computed ranks is evidence only after its matrices have been identified
with the specified code, action or quotient.

\begin{lemma}\label{ULCD.SquareZeroIncidence}
Let $R$ be a matrix over $F$ such that $RR^{\mathsf T}=0$. Then
$Q=R^{\mathsf T}R$ is symmetric, $Q^2=0$, and $\im Q$ is self-orthogonal.
If $Rj=0$ and $j^{\mathsf T}j=0$, then $\im Q+Fj$ is self-orthogonal too.
\end{lemma}
\begin{proof}
We have $Q^2=R^{\mathsf T}(RR^{\mathsf T})R=0$, and hence
$\langle Qx,Qy\rangle=x^{\mathsf T}Q^2y=0$.
Also $Qj=R^{\mathsf T}Rj=0$, so $j$ is orthogonal to $\im Q$.
The remaining pairing is the assumed norm of $j$.
\end{proof}

\begin{theorem}\label{ULCD.MatrixAlgebra}
Let $d\geq1$, and let $V=F^d$ carry a representation of a group generated by
$g_1,\ldots,g_s$, with matrices $M_1,\ldots,M_s$.
If $d^2$ products of these matrices, including the empty product if needed,
are linearly independent in $\operatorname{End}_F(V)$, then $V$ is
absolutely irreducible.
It suffices instead to exhibit a rank-one matrix $P=vu$ in the generated
unital algebra such that the columns $Mv$ and rows $uN$, as $M,N$ run
over products of the generators, each span a space of dimension $d$.
\end{theorem}
\begin{proof}
The first hypothesis says that the generated algebra is
$\operatorname{End}_F(V)$. Every group-invariant subspace is stable under
this algebra. If it contains a nonzero vector, an endomorphism can send
that vector to any specified vector of $V$, so the subspace is $V$.
Linear independence of the $d^2$ matrices persists under every field
extension: a nonzero full-size minor stays nonzero. The same argument
therefore proves irreducibility after extension.

For the second test, select products $M_i,N_j$ so that the $d$ columns
$M_iv$ and the $d$ rows $uN_j$ are bases. Their $d^2$ outer products
form a basis of the matrix space: changing either chosen basis to the
standard basis is invertible. Each outer product is $M_iPN_j$ and belongs
to the generated algebra. The first argument applies.
\end{proof}

\begin{theorem}\label{ULCD.Socle}
Let $V$ be a finite-dimensional $FG$-module, with an invariant filtration
$0=V_0<V_1<\cdots<V_s=V$ whose factors $D_i=V_i/V_{i-1}$ are simple.
Every simple submodule of $V$ is the image of an injective homomorphism
from one of the $D_i$. Consequently, if a semisimple submodule $N$ contains
the image of every homomorphism $D_i\to V$, then $N=\operatorname{Soc}(V)$.
For this containment it suffices to check every map in a vector-space
basis of each $\operatorname{Hom}_G(D_i,V)$.
\end{theorem}
\begin{proof}
Let $S\subseteq V$ be simple, and choose the least $i$ with $S\subseteq V_i$.
The invariant subspace $S\cap V_{i-1}$ is either zero or $S$; minimality
excludes the latter. Projection therefore embeds $S$ in the simple module
$D_i$, and is an isomorphism. Its inverse followed by inclusion in $V$
is the required homomorphism. Under the containment hypothesis every
simple submodule lies in $N$, so their sum lies in $N$. Conversely a
semisimple submodule is a sum of simple submodules and lies in that sum.
Finally an arbitrary homomorphism is a linear combination of basis maps;
its image lies in the sum of their images.
\end{proof}

In row coordinates the homomorphisms are exactly the solutions of
$XP_g=D_gX$ for each chosen generator. Solving the whole homogeneous
system provides the required basis. Checking a selected map would not
give the conclusion. Invariant prefixes are checked directly on each
generator; simplicity of their factors can be certified by
Theorem~\ref{ULCD.MatrixAlgebra}. Thus the criterion does not require a
classification of all irreducible representations of $G$.

\begin{theorem}\label{ULCD.Refinement}
Let $\mathcal S$ be a collection of distinct subsets of a finite set
$\Omega$. Form its bipartite incidence graph, with the two parts initially
coloured differently. Refine a partition by recording, for every vertex,
its old cell and its numbers of neighbours in every old cell, until no
cell splits. Choose $x_1,\ldots,x_t\in\Omega$ successively, separating
$x_i$ from its current cell and refining again. If that cell has size $b_i$
before separation, and all vertices of $\Omega$ have different colours
at the end, then
\[
 |\operatorname{Aut}(\Omega,\mathcal S)|\leq\prod_{i=1}^t b_i.
\]
In particular this bounds the coordinate-permutation group of any code
whose complete set of minimum supports is $\mathcal S$.
\end{theorem}
\begin{proof}
An incidence automorphism preserves the initial parts. Whenever it
preserves every old cell it also preserves the neighbour counts defining
the refined cells. After separating $x_1,\ldots,x_{i-1}$, the stabilizer
of those points therefore preserves every current cell. Its orbit of
$x_i$ has size at most $b_i$. At the last stage this stabilizer fixes
every point of $\Omega$ and hence every subset in $\mathcal S$.
It is trivial. Repeated orbit--stabilizer gives the inequality.
A coordinate permutation preserving a code preserves word weights and
therefore its complete minimum-support collection, proving the last claim.
\end{proof}

To obtain an equality, the application supplies actual permutations
preserving the code and verifies that they generate a subgroup of order
$\prod b_i$. The upper bound concerns the full group, independently of
how those permutations were discovered. For a graph with adjacency matrix
$A$, every graph automorphism preserves $\im A$; thus the additional
inclusions $G\leq\operatorname{Aut}(\Gamma)\leq\operatorname{Aut}(C_A)$
identify the full graph group whenever both outer groups have that order.

\section{Weights and dual weights}

The averaging theorem below limits a distance search when the code has
a transitive permutation group. A complete distribution requires either
enumeration of the whole code or another proof. In particular a count of
messages of bounded weight is not a whole-code weight distribution.

\begin{theorem}[MacWilliams identity]\label{ULCD.MacWilliams}
Let $C\leq\mathbb F_q^n$ have dimension $k$, and let $A_i$ be its number
of words of weight $i$. Write $A_j^\perp$ for the corresponding numbers
in the orthogonal dual. Then
\[
 A_j^\perp=q^{-k}\sum_{i=0}^n A_i
 \sum_{s=0}^j(-1)^s(q-1)^{j-s}
                 \binom is\binom{n-i}{j-s}.
\]
\end{theorem}
\begin{proof}
Choose a nontrivial additive character
$\chi:\mathbb F_q\to\mathbb C^\times$. Such a character is obtained
from a nonzero $\mathbb F_p$-linear functional to $\mathbb F_p$ and a
primitive $p$th root of unity, where $p$ is the characteristic.
For a nontrivial character on a finite additive group its sum is zero:
translation by an element on which it is not one multiplies the sum by
that value and also leaves it unchanged.
For fixed $x\in\mathbb F_q^n$, the map $c\mapsto c\cdot x$ is either
zero on $C$ or is onto $\mathbb F_q$, with equally sized fibres.
It follows that
\[
 q^{-k}\sum_{c\in C}\chi(c\cdot x)=\mathbf1_{x\in C^\perp}.
\]
Multiply by $z^{\mathrm{wt}(x)}$ and sum over all $x$. The finite sums
may be interchanged. For a fixed $c$ of weight $i$, the coordinate factors
in the inner sum are $1+(q-1)z$ at its zero entries and $1-z$ at its
nonzero entries. Hence
\[
 \sum_j A_j^\perp z^j
 =q^{-k}\sum_i A_i(1-z)^i(1+(q-1)z)^{n-i}.
\]
Taking the coefficient of $z^j$ proves the formula.
\end{proof}

This is a classical identity, included here to make its repeated use
explicit. For a matrix with $k$ independent rows, the map from its
$q^k$ coefficient vectors to its row space is bijective. Enumerating
each coefficient vector once therefore gives the complete $A_i$.
The displayed identity then determines the dual distribution. Integrality,
nonnegativity and totals provide consistency checks; these checks alone
would not prove that a proposed input distribution belongs to the code.
